\chapter{Corrections to v1.0.0}
\label{ch:corrections}

\texttt{v1.0.0} made claims that were wrong. This chapter says what each one was, what
was wrong with it, and what is true now. Values marked as \texttt{v1.0.0} values describe
the past and are not regenerated.

\section{Interior density}

\texttt{v1.0.0} reported an interior density error of 4.20 percent against a 5 percent
limit. That gate read densities in cell sorted order against positions in lattice order,
so it measured mismatched particles. With the index fixed, the \texttt{v1.0.0} mass gives
7.61 percent, which fails that limit. The explanation \texttt{v1.0.0} gave, a finite
smoothing radius on a jittered lattice, was also wrong: almost all of the residual was the
mass, which did not match the lattice laid out.

That mass started the dam break interior near 76 kPa against a hydrostatic 3.9 kPa, so the
\texttt{v1.0.0} fluid started with a pressure pulse. Maximum speed after 100 steps was
3.695 m/s and is now 0.625 m/s, still above the 0.098 m/s of free fall, so the start is
better but not quiet. The error is now \SphDensityErrorPercent{} percent by calibration,
with the limits stated in Section~\ref{sec:gates}.

\section{Fusion}

\texttt{v1.0.0} claimed a gain of 0.1 percent at its largest size, from a difference of
0.12 percent measured in separate blocks, with a session spread of about 3 percent. That
difference could not be told from noise. It also described the saving as a round trip of
the acceleration array, $24N$ bytes written and read, when the fused kernel still wrote
that array, and called the resulting 0.019 percent an absolute ceiling, which the
measurements in Section~\ref{sec:fusion} now exceed. It said fusion halves the memory
footprint of a step, which it does not: the fused path adds a second position buffer.

A first correction during this release said the saving was $12N$ bytes. That is also
incomplete, because the split path reads the acceleration array in both integrator
launches. Section~\ref{sec:fusion} gives the traffic as the code has it now, and the
paired ratios: fused is measurably faster up to one hundred and fifty thousand particles,
and the difference is not resolved at \LadderMaxN{}, for a reason not yet understood.

\section{Tile width}

\texttt{v1.0.0} defaulted to a tile width of 32, and the default is now \BestTileWidth{}.
Its tile sweep rows timed the force kernel alone for tiled, with particles that never
moved, against whole steps for fused, so any comparison across variants in that table was
void. Both variants now time a whole step.

\section{Provenance}

Both commit hashes recorded in the \texttt{v1.0.0} results do not exist in the published
history, because the history was rewritten before release. Every row now carries a hash
that resolves, here \texttt{\BenchCommit}, and a dirty flag, and
\texttt{scripts/check\_provenance.py} refuses a row that fails either.

\section{Energy gate}

The two body energy figure is dominated by single precision rounding, $9.9 \times 10^{-13}$
of truncation against \EnergyDriftMeasured{} measured, so it says nothing about the
integrator's order. \texttt{v1.0.0} presented it as evidence about the integrator. Order is
now tested in double precision.

\section{Numbers and prose that had drifted}

The Markdown documents quoted values that disagreed with the recorded results, among them
a streaming bandwidth of 545.8 GB/s and a throughput of 17.47 TFLOP/s, both \texttt{v1.0.0}
values that no longer match any file. They now carry marks that a script checks against
the recorded results. In the reports, \texttt{v1.0.0} said the tiling advantage grows with
particle count, which its own table contradicted; said the debug report has eleven entries
when it had fourteen; and called a tile width of 32 the slowest of the four measured when
its sweep had 64 slower. The benchmark protocol also changed, so timings here are not
directly comparable with \texttt{v1.0.0}.
